We have implemented the RC optimizations described in the previous
sections in the new compiler for the Lean programming language.
We have implemented all optimizations in Lean, and they are available
online~\footnote{\url{https://github.com/leanprover/lean4/tree/master/library/init/lean/compiler/ir}}.
At the time of writing, the compiler supports only one backend where we emit
C code. We are currently working on an LLVM backend for
our compiler. To test the efficiency of the compiler and RC
optimizations, we have devised a number of benchmarks\footnote{\url{https://github.com/leanprover/lean4/tree/IFL19/tests/bench}}
that aim to replicate common tasks performed in compilers and proof assistants.
All timings are arithmetic means of 50 runs as reported by the \texttt{temci}
benchmarking tool~\cite{bechberger16bachelorarbeit},
executed on a PC with an i7-3770 Intel CPU and 16~GB RAM running Ubuntu 18.04,
using Clang 9.0.0 for compiling the Lean runtime library
as well as the C code emitted by the Lean compiler.
\begin{itemize}
\item \verb!deriv! and \verb!const_fold! implement
  differentiation and constant folding, respectively, as examples of symbolic
  term manipulation where big expressions are constructed and transformed.
  We claim they reflect operations frequently performed by proof automation procedures used in theorem provers.
\item \verb!rbmap! stress tests the red-black tree implementation from the Lean standard library.
  The benchmarks \verb!rbmap_10! and \verb!rbmap_1! are two variants where
  we perform updates on shared trees.
\item \verb!parser! is the new parser we are developing for the next version of
  Lean. It is written purely in Lean (approximately $2000$ lines of code).
\item \verb!qsort! it is the basic quicksort algorithm for sorting arrays.
\item \verb!binarytrees! is taken from the Computer Languages Benchmarks
Game\footnote{\url{https://benchmarksgame-team.pages.debian.net/benchmarksgame/performance/binarytrees.html}}.
This benchmark is a simple adaption of Hans Boehm's \texttt{GCBench} benchmark\footnote{\url{http://hboehm.info/gc/gc_bench/}}.
The Lean version is a translation of the fastest, parallelized Haskell solution,
using \texttt{Task} in place of the Haskell \texttt{parallel} API.
\item \verb!unionfind! implements the union-find algorithm which is frequently used
  to implement decision procedures in automated reasoning. We use arrays to store
  the \emph{find} table, and thread the state using a state monad transformer
\end{itemize}

We have tested the impact of each optimization by selectively disabling it and
comparing the resulting runtime with the base runtime~(\cref{fig:leanbench}):
\begin{itemize}
\item \emph{-reuse} disables the insertion of \kw{reset}/\kw{reuse} operations
\item \emph{-borrow} disables borrow inference, assuming that all parameters are
  owned. Note that the compiler must still honor borrow annotations on builtins,
  which are unaffected.
\item \emph{-ST} uses atomic RC operations for all values
%% \item \emph{-MT} disables multi-threading and uses non-atomic RC operations for all values
%% \item ...
\end{itemize}

\begin{figure}
  \centering
  \begin{tabular}{lrrrr} \toprule
    & base & \emph{-reuse} & \emph{-borrow} & \emph{-ST} \\ \midrule
    \verb!binarytrees! & 1.0\ensuremath{\tilde{0}} & 0.9\ensuremath{\tilde{8}} & 1.1\ensuremath{\tilde{4}} & 1.2\ensuremath{\tilde{2}}\\
\verb!deriv! & 1.0\ensuremath{\tilde{0}} & 1.00 & 1.1\ensuremath{\tilde{6}} & 1.42\\
\verb!const_fold! & 1.00 & 1.64 & 0.90 & 1.23\\
\verb!parser! & 1.00 & 1.00 & 1.0\ensuremath{\tilde{0}} & 1.68\\
\verb!qsort! & 1.00 & 1.00 & 1.00 & 1.13\\
\verb!rbmap! & 1.0\ensuremath{\tilde{0}} & 3.23 & 1.0\ensuremath{\tilde{7}} & 1.71\\
\verb!  rbmap_10! & 1.49 & 3.62 & 1.52 & 2.43\\
\verb!  rbmap_1! & 4.7\ensuremath{\tilde{2}} & 5.42 & 4.47 & 8.02\\
\verb!unionfind! & 1.0\ensuremath{\tilde{0}} & 1.4\ensuremath{\tilde{1}} & 1.0\ensuremath{\tilde{0}} & 2.31\\
\midrule geom.\ mean & 1.24 & 1.74 & 1.2\ensuremath{\tilde{7}} & 1.89\\

    \bottomrule
  \end{tabular}
  \caption{Lean variant benchmarks, normalized by the base run time (\texttt{rbmap} for \texttt{rbmap\_*}). Digits whose
    order of magnitude is no larger than that of twice the standard deviation are marked by squiggly lines.}
  \label{fig:leanbench}
\end{figure}

The results show that the new \kw{reset} and \kw{reuse} instructions significantly improve
performance in the benchmarks \verb!const_fold!, \verb!rbmap!, and
\verb!unionfind!. The borrowed inference heuristic
provides significant speedups in benchmarks \verb!binarytrees! and
\verb!deriv!.
%% is also effective
%% in all but the \verb!const_fold! benchmark. Our thread safe RC scheme
%% improves performance in all benchmarks.
%% Finally, we have that \verb!unionfind1! is significantly
%% faster than \verb!unionfind2! demonstrating that our alternative definition for the state monad
%% (\cref{sec:reuse}) does indeed take advantage of the \kw{reset} and \kw{reuse} instructions.

We have also directly translated some of these programs to other statically
typed, functional languages: Haskell, OCaml, and Standard ML (\cref{fig:crossbench}).
For the latter we selected the compilers MLton~\cite{Weeks:2006:WCM:1159876.1159877},
which performs whole program optimization and can switch between multiple GC schemes at runtime,
and MLKit, which combines Region Inference and garbage collection~\cite{het02}.
While not primarily a functional language, we have also included Swift as a
popular statically typed language using reference counting.
For \verb!binarytrees!, we have used the original files and compiler flags from the fastest Benchmark
Game implementations. For Swift, we used the second-fastest, safe
implementation, which is much more comparable to the other versions than the
fastest one completely depending on unsafe code. The Benchmark Game does not include an SML version.
For \verb!qsort!, the Lean code is pure and relies on the fact that array
updates are destructive if the array is not shared. The Swift code behaves
similarly because Swift arrays are copy-on-write.
All other versions use destructive updates, using the \textit{ST} monad in the case of Haskell.

\begin{figure*}
  \centering
  \begin{tabular}{lrrrrrrrrrrrrrrrrrr} \toprule
    &\multicolumn{3}{c}{Lean 4} & \multicolumn{3}{c}{GHC 8.8.3} & \multicolumn{3}{c}{ocamlopt 4.10} & \multicolumn{3}{c}{MLton 20180207} & \multicolumn{3}{c}{MLKit 4.4.2} & \multicolumn{3}{c}{Swift 5.1.1} \\
    \cmidrule(lr{.5em}){2-4} \cmidrule(lr{.5em}){5-7} \cmidrule(lr{.5em}){8-10} \cmidrule(lr{.5em}){11-13} \cmidrule(lr{.5em}){14-16} \cmidrule(lr{.5em}){17-19}
    & Time & Del & CM & Time & GC & CM & Time & GC & CM & Time & GC & CM & Time & GC & CM & Time & GC & CM \\ \midrule
    \verb!binarytrees! & 1.0\ensuremath{\tilde{0}} & 38\% & 8 & 3.\ensuremath{\tilde{3}\tilde{8}} & 74\% & 13 & 1.3\ensuremath{\tilde{1}} & --- & 17 & --- & --- & --- & --- & --- & --- & 5.3\ensuremath{\tilde{5}} & 59\% & 10\\
\verb!deriv! & 1.0\ensuremath{\tilde{0}} & 2\ensuremath{\tilde{2}}\% & 12 & 2.23 & 43\% & 6 & 1.51 & 78\% & 6 & 0.98 & 21\% & 18 & 4.22 & 59\% & 20 & 3.7\ensuremath{\tilde{3}} & 42\% & 6\\
\verb!const_fold! & 1.00 & 1\ensuremath{\tilde{3}}\% & 18 & 2.73 & 59\% & 7 & 5.11 & 90\% & 5 & 1.15 & 29\% & 30 & 4.5\ensuremath{\tilde{9}} & 64\% & 14 & 6.3\ensuremath{\tilde{0}} & 49\% & 11\\
\verb!qsort! & 1.00 & 1\ensuremath{\tilde{0}}\% & 0 & 1.76 & 1\% & 0 & 1.37 & 34\% & 0 & 0.54 & \ensuremath{\tilde{0}}\% & 0 & 3.17 & \ensuremath{\tilde{0}}\% & 0 & 0.6\ensuremath{\tilde{4}} & \ensuremath{\tilde{0}}\% & 0\\
\verb!rbmap! & 1.0\ensuremath{\tilde{0}} & \ensuremath{\tilde{3}}\% & 3 & 2.44 & 36\% & 5 & 1.0\ensuremath{\tilde{3}} & 32\% & 5 & 3.25 & 30\% & 32 & 6.49 & 60\% & 10 & 8.0\ensuremath{\tilde{8}} & 59\% & 1\\
\verb!  rbmap_10! & 1.49 & 1\ensuremath{\tilde{3}}\% & 11 & 3.88 & 58\% & 7 & 1.9\ensuremath{\tilde{1}} & 59\% & 9 & 3.48 & 29\% & 32 & 9.34 & 72\% & 16 & 8.6\ensuremath{\tilde{7}} & 55\% & 3\\
\verb!  rbmap_1! & 4.7\ensuremath{\tilde{2}} & 28\% & 22 & 14.66 & 87\% & 8 & 9.20 & 89\% & 13 & 4.43 & 37\% & 30 & 15.50 & 84\% & 32 & 12.7\ensuremath{\tilde{6}} & 48\% & 14\\

    \bottomrule
  \end{tabular}
  \footnotetext{Because OCaml does not currently support multi-threading, the
    program uses the Unix \texttt{fork} call instead}

  \caption{Cross-language benchmarks. The measurements include wall clock time
    (normalized by the Lean base run time), GC time (in percent, as reported
    by the respective compiler), and last-level cache misses (CM, in million per
    second, as reported by \texttt{perf stat}). For Swift, we measure time spent
    in inc, dec, and deallocation runtime functions as GC time using \texttt{perf}.
    For Lean, the former are always inlined, so we can only measure object deletion
    time.}
  \label{fig:crossbench}
\end{figure*}
While the absolute runtimes in \cref{fig:crossbench}
are influenced by many factors other than the implementation
of garbage collection that make direct comparisons difficult,
the results still signify that both our garbage collection
and the overall runtime and compiler implementation are very
competitive.  We initially conjectured the good performance was a result of
reduced cache misses due to reusing allocations and a lack
of GC tracing.
However, the results demonstrate this is not the case.  The
only benchmark where the code generated by our compiler produces
significantly fewer cache misses is \verb!rbmap!.  Note that Lean is
5x as fast as OCaml on \verb!const_fold!  even though it
triggers more cache misses per second.  The results
suggest that Lean code is often faster in the benchmarks where
the code generated by other compilers spends a significant amount of time performing GC. Using
\verb!const_fold! as an example again, Lean spends only $17\%$
of the runtime deallocating memory, while OCaml spends $90\%$ in the
GC.  This comparison is not entirely precise since it does not include
the amount of time Lean spends updating reference counts, but it seems
to be the most plausible explanation for the difference in
performance.  The results for \verb!qsort! are surprising, the
Lean and Swift implementations outperforms all destructive ones but
MLton. We remark that MLton and Swift have a clear advantage since they use
arrays of unboxed machine integers, while Lean and the other compilers
use boxed values. We did not find a way to disable this optimization
in MLton or Swift to confirm our conjecture.  We believe this benchmark
demonstrates that our compiler allows programmers to write efficient
pure code that uses arrays and hashtables. For \verb!rbmap!, Lean is
much faster than all other systems except for OCaml. We imagined this would only be the case
when the tree was not shared. Then we devised the two variants
\verb!rbmap_10! and \verb!rbmap_1! which save the
current tree in a list after every tenth or every insertion, respectively.
The idea is to simulate the behavior of a backtracking search where we
store a copy of the state before each case-split. As expected, Lean's
performance decreases on these two variants since the tree is now a
shared value, and the time spent deallocating objects increases
substantially.  However, Lean still outperforms all systems but MLton
on \verb!rbmap_1!.  In all other systems but MLton and Swift, the time
spent on GC increases considerably. Finally, we point out that
MLton spends significantly less time on GC than the other languages using
a tracing GC in general.

% , and we should not
% discard reference counting as a viable memory management technique.

%%% Local Variables:
%%% mode: latex
%%% TeX-master: "refcount"
%%% End:
